\chapter{Ensemble Learning}

\begin{objectives}
After studying this chapter you will be able to:
\begin{itemize}
  \item \textbf{define} ensemble learning and explain \emph{why} combining models improves accuracy;
  \item compute the accuracy of a \textbf{majority vote} and explain the role of \textbf{diversity};
  \item explain \textbf{hard} and \textbf{soft voting} classifiers and work out an example;
  \item explain \textbf{bagging} and \textbf{pasting}, the \textbf{bootstrap} and \textbf{out-of-bag (OOB)} evaluation;
  \item describe how a \textbf{Random Forest} works and read \textbf{feature importances};
  \item run \textbf{AdaBoost} by hand and explain \textbf{gradient boosting};
  \item explain \textbf{stacking} with a meta-learner;
  \item compare bagging, boosting and stacking in terms of \textbf{bias and variance};
  \item answer 2-, 5- and 10-mark university questions on Unit~III (ensembles).
\end{itemize}
\end{objectives}

\begin{pullquote}[--- Francis Galton, 1907 (the ``wisdom of crowds'')]
At a county fair, 800 people guessed the weight of an ox. No one guessed it exactly, but the \emph{average} of all guesses was off by less than one percent.
\end{pullquote}

In Chapters~3 and~4 we studied \emph{single} models. This chapter asks a simple question:

\begin{center}
\emph{If one model is good, are many models better?}
\end{center}

The answer is usually \textbf{yes}, provided the models are \emph{different from each other}. Ensembles are behind most winning solutions in machine-learning competitions and many real systems.

%====================================================================
\section{Introduction to Ensemble Learning}

\subsection{The Idea: Wisdom of the Crowd}

Think of a \textbf{hospital}. For a difficult case, one doctor might make a mistake, but if five independent specialists examine the patient and the \emph{majority} agree on a diagnosis, you trust it much more. Or think of a \textbf{jury}: twelve people decide together rather than one judge alone.

\begin{definition}[: Ensemble learning]
\textbf{Ensemble learning} combines the predictions of \emph{several models} (called \textbf{base learners}, \textbf{weak learners} or \textbf{predictors}) to produce one final prediction that is usually \emph{more accurate and more stable} than the prediction of any single model.
\end{definition}

\begin{definition}[: Weak learner and strong learner]
A \textbf{weak learner} is a model only slightly better than random guessing (for two classes, accuracy a little above $50\%$), for example a very shallow decision tree. A \textbf{strong learner} is a model with high accuracy. A key theoretical result says that many weak learners can be combined into a strong learner.
\end{definition}

\subsection{Why Does It Work? A Simple Calculation}

Suppose we have $n$ classifiers, each correct with probability $p>0.5$, and \emph{assume their mistakes are independent}. If we take the \textbf{majority vote}, the ensemble is correct whenever more than half of the classifiers are correct. For $n$ classifiers (odd $n$) this probability is given by the binomial distribution:
\[
  P(\text{ensemble correct})=\sum_{k=\lceil (n+1)/2\rceil}^{n}\binom{n}{k}\,p^{k}(1-p)^{n-k}.
\]

\begin{example}[: Three classifiers, each 70\% accurate]
The majority is correct when at least $2$ of the $3$ are correct.

\medskip
\begin{center}
\begin{tabular}{@{}lcc@{}}
\toprule
\textbf{Number correct} & \textbf{Probability} & \textbf{Value}\\
\midrule
All 3 correct & $0.7^3$ & $0.343$\\
Exactly 2 correct & $3\times0.7^2\times0.3$ & $0.441$\\
\midrule
\textbf{At least 2 correct} & & $\mathbf{0.784}$\\
\bottomrule
\end{tabular}
\end{center}

\medskip
So three $70\%$ classifiers give a $78.4\%$ ensemble. With \textbf{five} classifiers:
\[
  P=\binom53 0.7^3 0.3^2+\binom54 0.7^4 0.3+0.7^5=0.3087+0.3602+0.1681=\mathbf{0.837}.
\]
\end{example}

\begin{figure}[!htbp]
\centering
\begin{tikzpicture}
\begin{axis}[width=11cm,height=6.6cm,xlabel={Number of independent classifiers $n$},ylabel={Majority-vote accuracy},
  xmin=0,xmax=52,ymin=0.5,ymax=1.02,grid=both,grid style={gray!18},
  legend style={at={(0.98,0.04)},anchor=south east,font=\small}]
\addplot[thick,mlred,smooth] coordinates {(1,0.55)(3,0.5748)(5,0.5931)(7,0.6083)(9,0.6214)(11,0.6331)(13,0.6437)(15,0.6535)(17,0.6626)(19,0.6710)(21,0.6790)(23,0.6865)(25,0.6937)(27,0.7005)(29,0.7070)(31,0.7132)(33,0.7191)(35,0.7249)(37,0.7304)(39,0.7357)(41,0.7408)(43,0.7458)(45,0.7506)(47,0.7552)(49,0.7597)(51,0.7641)};
\addlegendentry{each $p=0.55$}
\addplot[thick,mlorange,smooth] coordinates {(1,0.6)(3,0.648)(5,0.6826)(7,0.7102)(9,0.7334)(11,0.7535)(13,0.7712)(15,0.7869)(17,0.8011)(19,0.8139)(21,0.8256)(23,0.8364)(25,0.8462)(27,0.8553)(29,0.8638)(31,0.8716)(33,0.8789)(35,0.8857)(37,0.8920)(39,0.8979)(41,0.9035)(43,0.9087)(45,0.9135)(47,0.9181)(49,0.9224)(51,0.9265)};
\addlegendentry{each $p=0.60$}
\addplot[thick,mlgreen,smooth] coordinates {(1,0.7)(3,0.784)(5,0.8369)(7,0.8740)(9,0.9012)(11,0.9218)(13,0.9376)(15,0.9500)(17,0.9597)(19,0.9674)(21,0.9736)(23,0.9786)(25,0.9825)(27,0.9857)(29,0.9883)(31,0.9905)(33,0.9922)(35,0.9936)(37,0.9947)(39,0.9957)(41,0.9964)(43,0.9971)(45,0.9976)(47,0.9980)(49,0.9983)(51,0.9986)};
\addlegendentry{each $p=0.70$}
\end{axis}
\end{tikzpicture}
\caption{Accuracy of a majority vote of $n$ \emph{independent} classifiers, each with accuracy $p$.}\label{fig:majority}
\end{figure}

\begin{explain}
The horizontal axis is the \textbf{number of classifiers} voting and the vertical axis is the \textbf{accuracy of the majority vote}. Each curve is for a different accuracy of the \emph{individual} classifiers (red $55\%$, orange $60\%$, green $70\%$).

\medskip
Look at the \textbf{green curve}: with one classifier the accuracy is $70\%$; with $5$ it is $83.7\%$; with $21$ it is $97.4\%$; with $51$ it is almost $100\%$. Even the \textbf{red curve}, built from classifiers that are barely better than guessing ($55\%$), climbs steadily to about $76\%$ with $51$ voters.

\medskip
\textbf{Important warning:} this picture assumes the classifiers make \emph{independent} errors. In practice they are trained on the same data and their errors are \emph{correlated}, so real gains are smaller. This is why we work hard to make the members \textbf{diverse}.
\end{explain}

\begin{keyidea}
\textbf{Ensembles work best when the members are (i) each better than random guessing and (ii) as different from each other as possible.} Diversity can come from different \emph{algorithms}, different \emph{training samples}, different \emph{features}, or different \emph{hyperparameters}.
\end{keyidea}

\subsection{Types of Ensemble Methods}

\begin{figure}[!htbp]
\centering
\resizebox{\linewidth}{!}{%
\begin{tikzpicture}[
  n/.style={draw,rounded corners,thick,align=center,font=\small,minimum height=0.9cm,inner sep=5pt},
  arr/.style={thick}]
\node[n,fill=mlblue!25,draw=mlblue] (root) at (0,0) {\textbf{Ensemble methods}};
\node[n,fill=mlgreen!15,draw=mlgreen] (par) at (-5.4,-2.1) {\textbf{Parallel}\\(train independently)};
\node[n,fill=mlorange!15,draw=mlorange] (seq) at (0,-2.1) {\textbf{Sequential}\\(each fixes previous errors)};
\node[n,fill=mlpurple!12,draw=mlpurple] (stk) at (5.4,-2.1) {\textbf{Stacked}\\(learn how to combine)};
\node[n,fill=white] (v) at (-7.4,-4.3) {Voting};
\node[n,fill=white] (b) at (-5.4,-4.3) {Bagging\\Pasting};
\node[n,fill=white] (rf) at (-3.3,-4.3) {Random\\Forest};
\node[n,fill=white] (ad) at (-1.2,-4.3) {AdaBoost};
\node[n,fill=white] (gb) at (1.2,-4.3) {Gradient\\Boosting};
\node[n,fill=white] (st) at (5.4,-4.3) {Stacking\\(blending)};
\foreach \a in {par,seq,stk} \draw[arr] (root)--(\a);
\foreach \a in {v,b,rf} \draw[arr] (par)--(\a);
\foreach \a in {ad,gb} \draw[arr] (seq)--(\a);
\draw[arr] (stk)--(st);
\end{tikzpicture}}
\caption{Families of ensemble methods covered in this chapter.}\label{fig:enstax}
\end{figure}

\begin{explain}
Read the tree from top to bottom. The three coloured boxes are the three \textbf{strategies}.

\begin{itemize}
  \item \textbf{Parallel} (green): the models are trained \emph{independently} (possibly on different samples of the data) and their outputs are combined by voting or averaging. The aim is to reduce \textbf{variance}.
  \item \textbf{Sequential} (orange): models are trained \emph{one after another}, and each new model concentrates on the mistakes of the earlier ones. The aim is to reduce \textbf{bias}.
  \item \textbf{Stacked} (purple): a second-level model \emph{learns} the best way to combine the first-level models.
\end{itemize}
\end{explain}

\begin{theory}{Ensemble Learning}
\tpart{Definition}
\textbf{Ensemble learning} is a machine-learning technique that trains \emph{multiple base models} and \textbf{combines their predictions} (by voting, averaging or a learned combiner) to obtain better accuracy and robustness than any single model.

\tpart{Working}
\begin{enumerate}[itemsep=2pt,topsep=3pt]
  \item Create several base learners, making them \emph{diverse} (different algorithms, data samples or features).
  \item Train each base learner.
  \item Combine the predictions: majority/weighted vote for classification, average for regression, or a meta-model.
  \item Evaluate the ensemble on unseen test data.
\end{enumerate}

\tpart{Assumptions}
\begin{itemize}[itemsep=2pt,topsep=3pt]
  \item Each base learner is better than random guessing
  \item The learners make \emph{different} (ideally independent) errors
\end{itemize}

\tpart{Advantages}
\begin{itemize}[itemsep=2pt,topsep=3pt]
  \item Higher accuracy than a single model
  \item Reduces overfitting (lower variance) and is more robust to noise
  \item Works with almost any base algorithm
\end{itemize}

\tpart{Disadvantages}
\begin{itemize}[itemsep=2pt,topsep=3pt]
  \item More computation, memory and prediction time
  \item Harder to interpret than a single model
  \item Gains are small if the members are too similar
\end{itemize}

\tpart{Applications}
\begin{itemize}[itemsep=2pt,topsep=3pt]
  \item Credit scoring and fraud detection
  \item Medical diagnosis and recommender systems
  \item Winning solutions in data-science competitions
\end{itemize}
\end{theory}

%====================================================================
\section{Voting Classifiers}

\begin{theory}{Voting Classifier}
\tpart{Definition}
A \textbf{voting classifier} combines several \emph{different} classifiers trained on the \emph{same data}. In \textbf{hard voting} the final class is the one predicted by the \emph{majority} of the classifiers; in \textbf{soft voting} the predicted class \emph{probabilities} are averaged and the class with the highest average probability wins.

\tpart{Working}
\begin{enumerate}[itemsep=2pt,topsep=3pt]
  \item Train several diverse classifiers (e.g.\ logistic regression, SVM, random forest) on the same training set.
  \item \textbf{Hard voting:} each classifier predicts a class; count the votes; choose the most frequent class.
  \item \textbf{Soft voting:} each classifier outputs $P(\text{class})$; compute $\bar P(c)=\frac1m\sum_j P_j(c)$ (optionally weighted); choose $\argmax_c\bar P(c)$.
\end{enumerate}

\tpart{Assumptions}
\begin{itemize}[itemsep=2pt,topsep=3pt]
  \item Classifiers are reasonably accurate and diverse
  \item (Soft voting) every classifier can output well-calibrated probabilities
\end{itemize}

\tpart{Advantages}
\begin{itemize}[itemsep=2pt,topsep=3pt]
  \item Very simple to build and often beats the best single member
  \item Soft voting uses the classifiers' \emph{confidence}
\end{itemize}

\tpart{Disadvantages}
\begin{itemize}[itemsep=2pt,topsep=3pt]
  \item Poor members or highly similar members add little
  \item Needs several trained models, hence more cost
\end{itemize}

\tpart{Applications}
\begin{itemize}[itemsep=2pt,topsep=3pt]
  \item Combining different algorithms for text, image or tabular classification
\end{itemize}
\end{theory}

\begin{figure}[!htbp]
\centering
\begin{tikzpicture}[
  b/.style={draw,rounded corners,thick,align=center,font=\small,minimum height=0.9cm,minimum width=2.5cm},
  arr/.style={-{Stealth},thick}]
\node[b,fill=mlblue!12,draw=mlblue] (x) at (0,0) {New instance $\vect x$};
\node[b,fill=gray!12] (c1) at (4.4,1.6) {Classifier 1\\(logistic regression)};
\node[b,fill=gray!12] (c2) at (4.4,0) {Classifier 2\\(SVM)};
\node[b,fill=gray!12] (c3) at (4.4,-1.6) {Classifier 3\\(random forest)};
\node[b,fill=mlorange!14,draw=mlorange,minimum width=2.2cm] (v) at (9.2,0) {Combine\\(vote / average)};
\node[b,fill=mlgreen!18,draw=mlgreen,minimum width=2cm] (y) at (12.4,0) {Final\\prediction};
\foreach \c/\t in {c1/{A},c2/{B},c3/{B}} \draw[arr] (x)--(\c);
\draw[arr] (c1)--(v) node[midway,above,font=\footnotesize,sloped]{A};
\draw[arr] (c2)--(v) node[midway,above,font=\footnotesize]{B};
\draw[arr] (c3)--(v) node[midway,below,font=\footnotesize,sloped]{B};
\draw[arr] (v)--(y);
\end{tikzpicture}
\caption{A voting ensemble: three classifiers vote on the same instance.}\label{fig:voting}
\end{figure}

\begin{explain}
Follow the arrows from left to right. The \textbf{same instance} $\vect x$ is given to three \emph{different} classifiers. Each returns its own answer (labels A, B, B on the arrows). The orange box combines them: in \textbf{hard voting} it counts the votes (two B's against one A), so the \textbf{final prediction is B}. Notice that classifier~1 made a ``mistake'' but the ensemble still reached the answer the majority agreed on.
\end{explain}

\subsection{Hard vs.\ Soft Voting: Worked Example}

\begin{example}[: Which class is the new e-mail?]
Three classifiers estimate the probability that an e-mail is \textbf{spam} (class A); the other class is not-spam (B).

\medskip
\begin{center}
\begin{tabular}{@{}lccc@{}}
\toprule
\textbf{Classifier} & $P(\text{A})$ & $P(\text{B})$ & \textbf{Own prediction}\\
\midrule
1 (logistic regression) & $0.90$ & $0.10$ & A\\
2 (SVM) & $0.40$ & $0.60$ & B\\
3 (random forest) & $0.40$ & $0.60$ & B\\
\bottomrule
\end{tabular}
\end{center}

\medskip
\textbf{Hard voting:} votes are A, B, B $\Rightarrow$ \textbf{B} wins, $2$ to $1$.

\medskip
\textbf{Soft voting:} average probability of A is
\[
  \bar P(A)=\frac{0.90+0.40+0.40}{3}=\frac{1.70}{3}=0.567,\qquad
  \bar P(B)=0.433\ \Rightarrow\ \textbf{A wins}.
\]

\textbf{Why the difference?} Classifier~1 is \emph{very confident} ($90\%$) while the other two are only \emph{mildly} sure ($60\%$). Soft voting lets confident classifiers count for more.

\medskip
\textbf{Weighted soft voting:} if we trust classifier~1 three times as much (weights $3,1,1$):
\[
  \bar P(A)=\frac{3(0.90)+1(0.40)+1(0.40)}{3+1+1}=\frac{3.5}{5}=0.70\ \Rightarrow\ \textbf{A}.
\]
\end{example}

For \textbf{regression}, the ensemble simply \textbf{averages} the predictions: $\hat y=\frac1m\sum_{j=1}^{m}\hat y_j$.

\begin{center}
\begin{tabular}{@{}p{3.1cm}p{5.4cm}p{5.4cm}@{}}
\toprule
 & \textbf{Hard voting} & \textbf{Soft voting}\\
\midrule
Uses & Predicted class labels & Predicted probabilities\\
Rule & Majority vote & Highest average probability\\
Needs probabilities? & No & Yes\\
Accuracy & Good & Usually slightly better\\
\bottomrule
\end{tabular}
\end{center}

%====================================================================
\section{Bagging and Pasting}

Voting uses \emph{different algorithms} on the same data. Another way to create diversity is to use the \textbf{same algorithm} but train each copy on a \textbf{different random subset of the training data}.

\begin{definition}[: Bagging and pasting]
Both methods train $B$ copies of the same base learner, each on a random sample of the training set, and combine them by voting (classification) or averaging (regression).
\begin{itemize}
  \item \textbf{Bagging} (\emph{bootstrap aggregating}): samples are drawn \textbf{with replacement}.
  \item \textbf{Pasting}: samples are drawn \textbf{without replacement}.
\end{itemize}
\end{definition}

\subsection{The Bootstrap}

\begin{definition}[: Bootstrap sample]
A \textbf{bootstrap sample} of a dataset of size $n$ is obtained by drawing $n$ instances \emph{at random with replacement}. Some instances appear several times and some do not appear at all.
\end{definition}

\begin{example}[: Drawing one bootstrap sample]
Original data: instances $\{1,2,3,4,5\}$ ($n=5$). One bootstrap sample might be
\[
  \{2,\ 2,\ 5,\ 1,\ 5\}.
\]
Instance 2 and 5 appear twice, instance 1 once, and instances \textbf{3 and 4 are missing}. These missing instances are the \textbf{out-of-bag (OOB)} instances for this sample; they can be used to test the model trained on the sample.
\end{example}

\paragraph{How many instances are left out?} The probability that a particular instance is \emph{not} chosen in one draw is $1-\frac1n$; over $n$ draws it is $(1-\frac1n)^n\to e^{-1}\approx0.368$ for large $n$. Hence

\begin{itemize}
  \item about \textbf{63.2\%} of the instances appear in each bootstrap sample (at least once);
  \item about \textbf{36.8\%} are \textbf{out-of-bag}.
\end{itemize}

(For $n=100$: $0.99^{100}=0.366$; a simulation with $n=1000$ gave $62.2\%$ unique instances.)

\subsection{Algorithm and Aggregation}

\begin{figure}[!htbp]
\centering
\resizebox{\linewidth}{!}{%
\begin{tikzpicture}[
  b/.style={draw,rounded corners,thick,align=center,font=\small,minimum height=0.85cm},
  arr/.style={-{Stealth},thick}]
\node[b,fill=mlblue!12,draw=mlblue,minimum width=2.6cm,minimum height=1.4cm] (d) at (0,0) {Training set\\$(n$ instances$)$};
\node[b,fill=mlorange!14,draw=mlorange,minimum width=2.8cm] (s1) at (4.2,1.8) {Bootstrap sample 1};
\node[b,fill=mlorange!14,draw=mlorange,minimum width=2.8cm] (s2) at (4.2,0) {Bootstrap sample 2};
\node[b,fill=mlorange!14,draw=mlorange,minimum width=2.8cm] (s3) at (4.2,-1.8) {Bootstrap sample $B$};
\node[b,fill=gray!12,minimum width=1.9cm] (m1) at (8.2,1.8) {Model 1};
\node[b,fill=gray!12,minimum width=1.9cm] (m2) at (8.2,0) {Model 2};
\node[b,fill=gray!12,minimum width=1.9cm] (m3) at (8.2,-1.8) {Model $B$};
\node[b,fill=mlpurple!14,draw=mlpurple,minimum width=2.4cm,minimum height=1.4cm] (a) at (11.9,0) {Aggregate\\vote / average};
\node[b,fill=mlgreen!18,draw=mlgreen] (y) at (15,0) {Prediction};
\foreach \s in {s1,s2,s3} \draw[arr] (d)--(\s);
\foreach \s/\m in {s1/m1,s2/m2,s3/m3} \draw[arr] (\s)--(\m);
\foreach \m in {m1,m2,m3} \draw[arr] (\m)--(a);
\draw[arr] (a)--(y);
\node[font=\small,mlorange!80!black] at (4.2,-0.9) {$\vdots$};
\node[font=\small] at (8.2,-0.9) {$\vdots$};
\end{tikzpicture}}
\caption{The bagging procedure.}\label{fig:bagging}
\end{figure}

\begin{explain}
Start at the blue box on the left. (1)~From the training set we draw $B$ different \textbf{bootstrap samples} (orange). (2)~We train one model on each sample (grey); because the samples differ, the models differ. (3)~For a new instance, every model makes a prediction and the purple box \textbf{aggregates} them: the \textbf{most frequent class} for classification or the \textbf{average} for regression. The models can be trained at the same time (in parallel).
\end{explain}

\paragraph{Why does bagging reduce variance?} Suppose each of $B$ models has prediction variance $\sigma^2$ and any two models have correlation $\rho$. The variance of the average is
\begin{equation}\label{eq:bagvar}
  \Var(\bar f)=\rho\,\sigma^2+\frac{1-\rho}{B}\,\sigma^2 .
\end{equation}

\begin{example}[: Variance of an average of models]
Let $\sigma^2=4$ and $B=10$.
\begin{itemize}
  \item If the models are \textbf{independent} ($\rho=0$): $\Var=\frac{4}{10}=0.4$ (ten times smaller).
  \item If they are \textbf{highly correlated} ($\rho=0.5$): $\Var=0.5(4)+\frac{0.5}{10}(4)=2+0.2=2.2$ (only about half).
  \item Even with $B\to\infty$ the variance cannot go below $\rho\sigma^2=2$.
\end{itemize}
\textbf{Lesson:} averaging helps most when the models are \emph{uncorrelated}; bagging tries to decorrelate them using different bootstrap samples. Bagging reduces \textbf{variance} but \emph{not bias}, so it works best with \textbf{low-bias, high-variance} learners such as deep decision trees.
\end{example}

\subsection{Out-of-Bag (OOB) Evaluation}

Each model never saw about $36.8\%$ of the training instances. We can use these OOB instances as a \emph{free validation set}: for each instance, predict using only the models that did \textbf{not} have it in their sample, and compute the accuracy. The \textbf{OOB score} is a good estimate of test accuracy \emph{without needing a separate validation set}.

\begin{theory}{Bagging and Pasting}
\tpart{Definition}
\textbf{Bagging} (bootstrap aggregating) trains many copies of the same learner, each on a \emph{bootstrap sample} (random sample with replacement) of the training data, and combines their outputs by voting or averaging. \textbf{Pasting} is the same but samples \emph{without replacement}.

\tpart{Working}
\begin{enumerate}[itemsep=2pt,topsep=3pt]
  \item Draw $B$ random samples from the training set (with replacement for bagging).
  \item Train one base learner (e.g.\ a deep decision tree) on each sample.
  \item For a new instance collect the $B$ predictions.
  \item Classification: majority vote; regression: average.
  \item Optionally estimate accuracy using the out-of-bag instances.
\end{enumerate}

\tpart{Assumptions}
\begin{itemize}[itemsep=2pt,topsep=3pt]
  \item The base learner is unstable (small data changes cause large changes in the model)
  \item Bootstrap samples are representative of the population
\end{itemize}

\tpart{Advantages}
\begin{itemize}[itemsep=2pt,topsep=3pt]
  \item Reduces variance and overfitting
  \item Parallel training; free OOB validation
  \item Works with any base learner
\end{itemize}

\tpart{Disadvantages}
\begin{itemize}[itemsep=2pt,topsep=3pt]
  \item Does not reduce bias
  \item Less interpretable and slower than one model
  \item Little gain for stable learners (e.g.\ linear regression)
\end{itemize}

\tpart{Applications}
\begin{itemize}[itemsep=2pt,topsep=3pt]
  \item Improving decision trees (Random Forests)
  \item Noisy data problems, regression and classification
\end{itemize}
\end{theory}

\begin{center}
\begin{tabular}{@{}p{3.2cm}p{5.4cm}p{5.4cm}@{}}
\toprule
 & \textbf{Bagging} & \textbf{Pasting}\\
\midrule
Sampling & With replacement (bootstrap) & Without replacement\\
Duplicates in a sample? & Yes & No\\
OOB evaluation? & Yes (about $37\%$ per model) & Not directly\\
Diversity & Slightly higher & Slightly lower\\
Typical bias / variance & Slightly higher bias, lower variance & Slightly lower bias\\
\bottomrule
\end{tabular}
\end{center}

%====================================================================
\section{Random Forests}

A \textbf{Random Forest} is the most popular ensemble method. It is \textbf{bagging of decision trees} with one extra trick that makes the trees even more different.

\begin{definition}[: Random Forest]
A \textbf{Random Forest} is an ensemble of many decision trees, each trained on a \emph{bootstrap sample} of the data, where at \emph{every split} the tree considers only a \textbf{random subset of the features} (typically $\sqrt p$ for classification and $p/3$ for regression, with $p$ the total number of features). The forest predicts by \emph{majority vote} (classification) or \emph{average} (regression).
\end{definition}

\paragraph{Why the random feature subset?} If one feature is very strong, \emph{every} bagged tree would split on it first and all trees would look alike (high correlation $\rho$ in Equation~\eqref{eq:bagvar}). Forcing each split to choose from a random subset \textbf{decorrelates} the trees, so averaging them reduces the variance much more.

\begin{figure}[!htbp]
\centering
\resizebox{\linewidth}{!}{%
\begin{tikzpicture}[
  b/.style={draw,rounded corners,thick,align=center,font=\small,minimum height=0.85cm},
  arr/.style={-{Stealth},thick},
  nd/.style={circle,draw,thick,fill=mlblue!15,draw=mlblue,minimum size=0.55cm,inner sep=0pt},
  lf/.style={rectangle,draw,thick,minimum size=0.42cm,inner sep=0pt}]
\node[b,fill=mlblue!12,draw=mlblue] (d) at (0,0) {Training data};
% tree 1
\node[nd] (a1) at (4.6,2.5) {}; \node[nd] (a2) at (4.0,1.5) {}; \node[nd] (a3) at (5.2,1.5) {};
\node[lf,fill=mlgreen!30] (a4) at (3.6,0.7) {}; \node[lf,fill=mlred!25] (a5) at (4.4,0.7) {};
\node[lf,fill=mlred!25] (a6) at (5.2,0.7) {};
\draw[thick] (a1)--(a2)--(a4); \draw[thick] (a2)--(a5); \draw[thick] (a1)--(a3)--(a6);
\node[font=\small] at (4.4,3.2) {Tree 1 (sample 1, features $\{1,3\}$)};
% tree 2
\node[nd] (b1) at (4.6,-0.7) {}; \node[nd] (b2) at (4.0,-1.7) {}; \node[nd] (b3) at (5.2,-1.7) {};
\node[lf,fill=mlgreen!30] (b4) at (3.6,-2.5) {}; \node[lf,fill=mlgreen!30] (b5) at (4.4,-2.5) {};
\node[lf,fill=mlred!25] (b6) at (5.2,-2.5) {};
\draw[thick] (b1)--(b2)--(b4); \draw[thick] (b2)--(b5); \draw[thick] (b1)--(b3)--(b6);
\node[font=\small] at (4.4,0.05) {Tree 2 (sample 2, features $\{2,4\}$)};
% tree 3
\node[nd] (c1) at (4.6,-3.9) {}; \node[nd] (c2) at (4.0,-4.9) {}; \node[nd] (c3) at (5.2,-4.9) {};
\node[lf,fill=mlgreen!30] (c4) at (3.6,-5.7) {}; \node[lf,fill=mlgreen!30] (c5) at (4.4,-5.7) {};
\node[lf,fill=mlgreen!30] (c6) at (5.2,-5.7) {};
\draw[thick] (c1)--(c2)--(c4); \draw[thick] (c2)--(c5); \draw[thick] (c1)--(c3)--(c6);
\node[font=\small] at (4.4,-3.2) {Tree 3 (sample 3, features $\{1,4\}$)};
\node[b,fill=mlgreen!20,draw=mlgreen] (p1) at (8.4,1.1) {Tree 1 votes\\\textbf{Red}};
\node[b,fill=mlgreen!20,draw=mlgreen] (p2) at (8.4,-1.7) {Tree 2 votes\\\textbf{Green}};
\node[b,fill=mlgreen!20,draw=mlgreen] (p3) at (8.4,-4.8) {Tree 3 votes\\\textbf{Green}};
\node[b,fill=mlpurple!14,draw=mlpurple,minimum height=1.4cm] (m) at (11.8,-1.9) {Majority vote\\$\Rightarrow$ \textbf{Green}};
\draw[arr] (d)--(3.2,0.6); \draw[arr] (d)--(3.2,-1.4); \draw[arr] (d) to[out=-60,in=180] (3.2,-4.8);
\draw[arr] (5.9,0.9)--(p1); \draw[arr] (5.9,-1.8)--(p2); \draw[arr] (5.9,-4.9)--(p3);
\draw[arr] (p1)--(m); \draw[arr] (p2)--(m); \draw[arr] (p3)--(m);
\end{tikzpicture}}
\caption{A Random Forest of three trees. Each tree sees a different bootstrap sample and different candidate features.}\label{fig:rf}
\end{figure}

\begin{explain}
The training data (left) is turned into \textbf{three different trees}. Each tree is trained on its \emph{own bootstrap sample} and, at each split, can only choose from \emph{its own random features} (shown in brackets), so the trees have different shapes and make different mistakes. The small squares at the bottom of each tree are \textbf{leaves}: red or green squares show the class that leaf predicts. For a new instance each tree sends it down to a leaf and \textbf{votes} for that leaf's class: Tree~1 says Red; Trees~2 and~3 say Green. The forest takes the \textbf{majority}: \textbf{Green}.
\end{explain}

\begin{example}[: Voting in a forest of five trees]
A forest has $5$ trees. For a new patient the trees predict: Disease, Healthy, Disease, Disease, Healthy.

\medskip
\begin{center}
\begin{tabular}{@{}lcc@{}}
\toprule
 & \textbf{Disease} & \textbf{Healthy}\\
\midrule
Votes & 3 & 2\\
Estimated probability & $3/5=0.60$ & $2/5=0.40$\\
\bottomrule
\end{tabular}
\end{center}
\medskip
The forest predicts \textbf{Disease} (with a ``confidence'' of $60\%$). For regression the five tree outputs would be averaged.
\end{example}

\subsection{Feature Importance}

A great by-product: a forest tells us \textbf{which features matter}. A feature's \textbf{importance} is the average total decrease in impurity (Gini or entropy) from all splits that use it, averaged over all trees and normalised so the importances sum to $1$.

\begin{figure}[!htbp]
\centering
\begin{tikzpicture}
\begin{axis}[xbar,width=10.5cm,height=5.4cm,xlabel={Feature importance},
  symbolic y coords={Sepal width,Sepal length,Petal length,Petal width},ytick=data,
  xmin=0,xmax=0.55,bar width=14pt,nodes near coords,nodes near coords align={horizontal},
  every node near coord/.append style={font=\small},xmajorgrids,grid style={gray!18}]
\addplot[fill=mlblue!60,draw=mlblue] coordinates {(0.030,Sepal width)(0.098,Sepal length)(0.436,Petal length)(0.437,Petal width)};
\end{axis}
\end{tikzpicture}
\caption{Feature importances of a Random Forest (200 trees) on the Iris dataset.}\label{fig:imp}
\end{figure}

\begin{explain}
Each bar is one of the four measurements of an iris flower; a \textbf{longer bar means the forest relied on that feature more}. The two \textbf{petal} measurements together account for about $87\%$ of the importance ($0.436+0.437$), while the \emph{sepal width} contributes only $3\%$. This agrees with the single tree of Chapter~3, which separated all three species using only the petal width. Such a ranking helps in \textbf{feature selection} and in explaining the model to non-experts.
\end{explain}

\begin{theory}{Random Forest}
\tpart{Definition}
A \textbf{Random Forest} is an ensemble of decision trees built by \emph{bagging}, in which each split considers only a \emph{random subset of features}; the final output is the majority vote (classification) or the mean (regression) of the trees.

\tpart{Working}
\begin{enumerate}[itemsep=2pt,topsep=3pt]
  \item Draw $B$ bootstrap samples from the training data.
  \item For each sample grow a decision tree; at \emph{each node} choose the best split from a random subset of $m\approx\sqrt p$ features.
  \item Grow the trees deep (little or no pruning).
  \item Predict by majority vote or average of all trees.
  \item Compute the OOB score and the feature importances.
\end{enumerate}

\tpart{Assumptions}
\begin{itemize}[itemsep=2pt,topsep=3pt]
  \item Individual trees are low-bias but high-variance
  \item Trees are weakly correlated because of the randomisation
\end{itemize}

\tpart{Advantages}
\begin{itemize}[itemsep=2pt,topsep=3pt]
  \item High accuracy; handles many features and mixed data types
  \item Much less overfitting than a single tree
  \item Gives feature importance and OOB error; needs little tuning
\end{itemize}

\tpart{Disadvantages}
\begin{itemize}[itemsep=2pt,topsep=3pt]
  \item Less interpretable than one tree; larger memory
  \item Slower prediction with many trees
  \item May be biased towards features with many categories
\end{itemize}

\tpart{Applications}
\begin{itemize}[itemsep=2pt,topsep=3pt]
  \item Credit-risk and churn prediction
  \item Medical diagnosis, bioinformatics, remote sensing
  \item Feature selection
\end{itemize}
\end{theory}

\paragraph{Main hyperparameters.}
\begin{itemize}
  \item \texttt{n\_estimators}: number of trees (more is better but slower; $100$--$500$ is typical).
  \item \texttt{max\_features}: features tried at each split (default $\sqrt p$).
  \item \texttt{max\_depth}, \texttt{min\_samples\_leaf}: limit the size of each tree.
\end{itemize}

\paragraph{Extra-Trees.} \emph{Extremely Randomised Trees} go one step further and also choose the \emph{split thresholds} at random: faster to train and even lower variance, at the price of a little more bias.

%====================================================================
\section{Boosting}

Bagging trains models \emph{in parallel} to reduce variance. \textbf{Boosting} trains them \textbf{one after another}, each model trying to correct the mistakes of the previous ones. It reduces \textbf{bias}.

\begin{definition}[: Boosting]
\textbf{Boosting} is a sequential ensemble method that combines many \emph{weak learners} into a strong learner by training them in sequence, where each new learner focuses on the instances that the earlier learners handled badly, and the final prediction is a \emph{weighted} combination of all learners.
\end{definition}

\begin{keyidea}
\textbf{Analogy.} A student takes a test, then studies \emph{only the questions answered wrongly}, then takes another test, and so on. After several rounds the student has covered the hard questions that a single reading would have missed.
\end{keyidea}

\subsection{AdaBoost (Adaptive Boosting)}

\textbf{AdaBoost} keeps a \textbf{weight} on every training instance. After each round it \textbf{increases the weights of misclassified instances} so the next learner pays more attention to them. Labels are coded $y_i\in\{-1,+1\}$.

\begin{definition}[: AdaBoost algorithm]
\begin{enumerate}
  \item Initialise weights $w_i=\dfrac1n$ for all $n$ instances.
  \item For $t=1,\dots,T$:
  \begin{enumerate}[label=(\alph*)]
    \item Train a weak learner $h_t$ on the \emph{weighted} data.
    \item Weighted error: $\varepsilon_t=\sum_{i:\,h_t(\vect x_i)\neq y_i}w_i$.
    \item Learner weight (``say''): $\alpha_t=\dfrac12\ln\dfrac{1-\varepsilon_t}{\varepsilon_t}$.
    \item Update weights: $w_i\leftarrow w_i\,e^{-\alpha_t\,y_i\,h_t(\vect x_i)}$, then divide by the sum so that $\sum w_i=1$.\\
    (Correct instances are multiplied by $e^{-\alpha_t}<1$, wrong ones by $e^{+\alpha_t}>1$.)
  \end{enumerate}
  \item Final classifier: $H(\vect x)=\operatorname{sign}\Big(\sum_{t=1}^{T}\alpha_t\,h_t(\vect x)\Big)$.
\end{enumerate}
\end{definition}

\paragraph{Meaning of $\alpha_t$.} A learner with small error gets a \textbf{large} weight, one with error $0.5$ (random guessing) gets $\alpha=0$, and one with error above $0.5$ would get a negative weight.

\medskip
\begin{center}
\begin{tabular}{@{}c|ccccc@{}}
\toprule
Error $\varepsilon_t$ & $0.05$ & $0.10$ & $0.20$ & $0.30$ & $0.50$\\
\midrule
$\alpha_t=\frac12\ln\frac{1-\varepsilon_t}{\varepsilon_t}$ & $1.47$ & $1.10$ & $0.69$ & $0.42$ & $0$\\
\bottomrule
\end{tabular}
\end{center}

\subsubsection*{End-to-end example: five instances, three rounds}

True labels $y=(+1,+1,-1,-1,+1)$ for instances $1,\dots,5$. We use weak learners that are each wrong on only one instance.

\begin{example}[: AdaBoost by hand]
\textbf{Round 1.} Weights $w=(0.2,0.2,0.2,0.2,0.2)$.

Weak learner $h_1$ predicts $(+,+,-,+,+)$: it is wrong only on \textbf{instance 4}.
\[
  \varepsilon_1=0.2,\qquad \alpha_1=\tfrac12\ln\tfrac{0.8}{0.2}=\tfrac12\ln4=0.693 .
\]
Update: correct instances $\times e^{-0.693}=\times0.5$; instance 4 $\times e^{0.693}=\times2$.
\[
  (0.1,\,0.1,\,0.1,\,0.4,\,0.1)\ \xrightarrow{\ \text{sum}=0.8\ }\ w=(0.125,\,0.125,\,0.125,\,\mathbf{0.5},\,0.125).
\]

\medskip
\textbf{Round 2.} Instance 4 now carries half of the total weight, so $h_2$ must get it right. Say $h_2$ predicts $(-,+,-,-,+)$: wrong only on \textbf{instance 1}.
\[
  \varepsilon_2=0.125,\qquad \alpha_2=\tfrac12\ln\tfrac{0.875}{0.125}=\tfrac12\ln7=0.973 .
\]
Correct $\times e^{-0.973}=0.378$; instance 1 $\times e^{0.973}=2.646$. After normalising:
\[
  w=(\mathbf{0.500},\,0.071,\,0.071,\,0.286,\,0.071).
\]

\medskip
\textbf{Round 3.} Now instance 1 dominates. Say $h_3$ predicts $(+,+,+,-,+)$: wrong only on \textbf{instance 3}.
\[
  \varepsilon_3=0.071,\qquad \alpha_3=\tfrac12\ln\tfrac{0.929}{0.071}=\tfrac12\ln13=1.282 .
\]

\medskip
\textbf{Final ensemble:} $H(\vect x)=\operatorname{sign}\big(0.693\,h_1+0.973\,h_2+1.282\,h_3\big)$.

\begin{center}
\begin{tabular}{@{}c|cccccc|c@{}}
\toprule
Instance & $y$ & $h_1$ & $h_2$ & $h_3$ & Score & $H$ & Correct?\\
\midrule
1 & $+$ & $+$ & $-$ & $+$ & $0.693-0.973+1.282=+1.00$ & $+$ & \checkmark\\
2 & $+$ & $+$ & $+$ & $+$ & $+2.95$ & $+$ & \checkmark\\
3 & $-$ & $-$ & $-$ & $+$ & $-0.693-0.973+1.282=-0.38$ & $-$ & \checkmark\\
4 & $-$ & $+$ & $-$ & $-$ & $0.693-0.973-1.282=-1.56$ & $-$ & \checkmark\\
5 & $+$ & $+$ & $+$ & $+$ & $+2.95$ & $+$ & \checkmark\\
\bottomrule
\end{tabular}
\end{center}

\textbf{Conclusion:} each weak learner made a mistake, but the weighted ensemble classifies \emph{all five} instances correctly.
\end{example}

\begin{figure}[!htbp]
\centering
\begin{tikzpicture}
\begin{axis}[ybar,width=12cm,height=6.2cm,bar width=9pt,xlabel={Instance},ylabel={Weight $w_i$},
  symbolic x coords={1,2,3,4,5},xtick=data,ymin=0,ymax=0.6,enlarge x limits=0.15,
  legend style={at={(0.5,1.04)},anchor=south,legend columns=3,font=\small},ymajorgrids,grid style={gray!18}]
\addplot[fill=mlblue!55,draw=mlblue] coordinates {(1,0.2)(2,0.2)(3,0.2)(4,0.2)(5,0.2)};
\addplot[fill=mlorange!65,draw=mlorange] coordinates {(1,0.125)(2,0.125)(3,0.125)(4,0.5)(5,0.125)};
\addplot[fill=mlgreen!60,draw=mlgreen] coordinates {(1,0.5)(2,0.071)(3,0.071)(4,0.286)(5,0.071)};
\legend{Round 1,Round 2,Round 3}
\end{axis}
\end{tikzpicture}
\caption{How AdaBoost's instance weights change from round to round in the worked example.}\label{fig:adaw}
\end{figure}

\begin{explain}
Each group of bars is one training instance; the three bars show its weight at the start of Round~1 (blue), Round~2 (orange) and Round~3 (green).

\begin{itemize}
  \item \textbf{Round 1 (blue):} all instances are equal ($0.2$).
  \item \textbf{Round 2 (orange):} instance~4, which learner~1 got wrong, has jumped to $0.5$.
  \item \textbf{Round 3 (green):} instance~1, which learner~2 got wrong, now has the biggest weight ($0.5$) while instance~4 (now handled) has dropped back to $0.286$.
\end{itemize}
The \textbf{spotlight moves} to whichever instances are still being misclassified. This is how boosting makes later learners specialise in the hard cases.
\end{explain}

\begin{theory}{AdaBoost}
\tpart{Definition}
\textbf{AdaBoost} (Adaptive Boosting) is a \emph{sequential} ensemble algorithm that builds a strong classifier as a \emph{weighted vote of weak learners} (usually decision stumps), re-weighting the training instances after every round so that misclassified instances get more attention.

\tpart{Working}
\begin{enumerate}[itemsep=2pt,topsep=3pt]
  \item Start with equal instance weights $w_i=1/n$.
  \item Train a weak learner on the weighted data and compute its weighted error $\varepsilon_t$.
  \item Compute its weight $\alpha_t=\frac12\ln\frac{1-\varepsilon_t}{\varepsilon_t}$.
  \item Increase the weights of misclassified instances, decrease the others, and normalise.
  \item Repeat for $T$ rounds; predict $H(\vect x)=\operatorname{sign}\sum_t\alpha_th_t(\vect x)$.
\end{enumerate}

\tpart{Assumptions}
\begin{itemize}[itemsep=2pt,topsep=3pt]
  \item Each weak learner is slightly better than random guessing ($\varepsilon_t<0.5$)
  \item Data are not dominated by label noise
\end{itemize}

\tpart{Advantages}
\begin{itemize}[itemsep=2pt,topsep=3pt]
  \item Turns weak learners into a strong one; reduces bias
  \item Few hyperparameters; simple to implement
\end{itemize}

\tpart{Disadvantages}
\begin{itemize}[itemsep=2pt,topsep=3pt]
  \item Sensitive to noisy data and outliers (they keep gaining weight)
  \item Sequential, so it cannot be parallelised
\end{itemize}

\tpart{Applications}
\begin{itemize}[itemsep=2pt,topsep=3pt]
  \item Face detection (Viola--Jones)
  \item Text and customer classification
\end{itemize}
\end{theory}

\subsection{Gradient Boosting}

\textbf{Gradient boosting} generalises the idea. Instead of re-weighting instances it fits each new model to the \textbf{residuals} (errors) of the current ensemble.

\begin{definition}[: Gradient boosting (regression with squared loss)]
\begin{enumerate}
  \item Start with a constant prediction $F_0(\vect x)=\bar y$ (the mean).
  \item For $m=1,\dots,M$:
  \begin{enumerate}[label=(\alph*)]
    \item Compute the \textbf{residuals} $r_i=y_i-F_{m-1}(\vect x_i)$ (they are the negative gradient of the squared-error loss).
    \item Fit a small tree $h_m$ to predict the residuals.
    \item Update: $F_m(\vect x)=F_{m-1}(\vect x)+\eta\,h_m(\vect x)$, where $\eta\in(0,1]$ is the \textbf{learning rate} (shrinkage).
  \end{enumerate}
  \item The final model is $F_M(\vect x)=\bar y+\eta\sum_{m=1}^{M}h_m(\vect x)$.
\end{enumerate}
\end{definition}

\begin{example}[: Gradient boosting on four points]
Predict a salary (in thousand rupees) from years of experience $x$:

\medskip
\begin{center}
\begin{tabular}{@{}c|cccc@{}}
\toprule
$x$ & 1 & 2 & 3 & 4\\
$y$ & 10 & 12 & 20 & 22\\
\bottomrule
\end{tabular}
\end{center}

\medskip
Use a decision stump that splits at $x\le2.5$ and $\eta=0.5$.

\medskip
\textbf{Step 0.} $F_0=\bar y=\frac{10+12+20+22}{4}=16$ for all $x$. Residuals $r=(-6,-4,4,6)$; $\text{MSE}=\frac{36+16+16+36}{4}=\mathbf{26}$.

\medskip
\textbf{Step 1.} The stump predicts the mean residual on each side: left $\frac{-6-4}{2}=-5$, right $\frac{4+6}{2}=5$.
\[
  F_1=F_0+0.5\,h_1=(16-2.5,\ 16-2.5,\ 16+2.5,\ 16+2.5)=(13.5,\ 13.5,\ 18.5,\ 18.5).
\]
New residuals $(10-13.5,\ 12-13.5,\ 20-18.5,\ 22-18.5)=(-3.5,-1.5,1.5,3.5)$; $\text{MSE}=\frac{12.25+2.25+2.25+12.25}{4}=\mathbf{7.25}$.

\medskip
\textbf{Step 2.} Stump on the new residuals: left $-2.5$, right $+2.5$.
\[
  F_2=F_1+0.5\,h_2=(12.25,\ 12.25,\ 19.75,\ 19.75).
\]
Residuals $(-2.25,-0.25,0.25,2.25)$; $\text{MSE}=\frac{5.06+0.06+0.06+5.06}{4}=\mathbf{2.56}$.

\medskip
\begin{center}
\begin{tabular}{@{}lccc@{}}
\toprule
\textbf{Stage} & $F_0$ & $F_1$ & $F_2$\\
\midrule
MSE & 26.0 & 7.25 & 2.56\\
\bottomrule
\end{tabular}
\end{center}
The error falls at every stage. Prediction for $x=3$: $F_2(3)=19.75$ (true value 20).
\end{example}

\begin{figure}[!htbp]
\centering
\begin{tikzpicture}
\begin{axis}[width=11cm,height=6.4cm,xlabel={Years of experience $x$},ylabel={Salary $y$},
  xmin=0.5,xmax=4.5,ymin=8,ymax=24,grid=both,grid style={gray!18},xtick={1,2,3,4},
  legend style={at={(0.02,0.98)},anchor=north west,font=\small}]
\addplot[only marks,mark=*,mlblue,mark size=3.2pt] coordinates {(1,10)(2,12)(3,20)(4,22)};
\addlegendentry{data $y$}
\addplot[thick,gray,dashed] coordinates {(0.5,16)(4.5,16)};
\addlegendentry{$F_0$ (mean)}
\addplot[thick,mlorange] coordinates {(0.5,13.5)(2.5,13.5)(2.5,18.5)(4.5,18.5)};
\addlegendentry{$F_1$}
\addplot[very thick,mlgreen] coordinates {(0.5,12.25)(2.5,12.25)(2.5,19.75)(4.5,19.75)};
\addlegendentry{$F_2$}
\end{axis}
\end{tikzpicture}
\caption{Gradient boosting: each stage moves the step-function prediction closer to the data.}\label{fig:gbfit}
\end{figure}

\begin{explain}
The blue dots are the data. The \textbf{grey dashed line} is the starting guess $F_0$: the mean, a flat line that is wrong for every point. The \textbf{orange step function} $F_1$ moves the left points down and the right points up, getting closer. The \textbf{green step function} $F_2$ is closer still: its left step ($12.25$) is near the left points ($10,12$) and its right step ($19.75$) is near ($20,22$). Each stage \emph{adds a small correction} to fix the remaining error. With many stages and deeper trees, the curve can follow complex patterns.
\end{explain}

\paragraph{Learning rate and number of stages.} A small $\eta$ (e.g.\ $0.1$) makes learning slower but usually generalises better; it must be balanced with the number of trees $M$ (small $\eta$ needs large $M$). Too many stages can \textbf{overfit}; use early stopping on a validation set. Popular fast implementations are \textbf{XGBoost}, \textbf{LightGBM} and \textbf{CatBoost}.

\begin{theory}{Gradient Boosting}
\tpart{Definition}
\textbf{Gradient boosting} builds an additive model $F_M(\vect x)=F_0+\eta\sum_mh_m(\vect x)$ \emph{stage by stage}, where each new weak learner $h_m$ is fitted to the \textbf{negative gradient of the loss} (the \emph{residuals}, for squared error) of the current model.

\tpart{Working}
\begin{enumerate}[itemsep=2pt,topsep=3pt]
  \item Initialise with a constant prediction (e.g.\ the mean).
  \item Compute the residuals of the current model.
  \item Fit a small decision tree to the residuals.
  \item Add the tree to the model scaled by the learning rate $\eta$.
  \item Repeat for $M$ stages (stop early on a validation set).
\end{enumerate}

\tpart{Assumptions}
\begin{itemize}[itemsep=2pt,topsep=3pt]
  \item A differentiable loss function suited to the task
  \item Weak learners are simple (shallow trees)
\end{itemize}

\tpart{Advantages}
\begin{itemize}[itemsep=2pt,topsep=3pt]
  \item Very high accuracy on structured/tabular data
  \item Works for regression, classification and ranking with any differentiable loss
\end{itemize}

\tpart{Disadvantages}
\begin{itemize}[itemsep=2pt,topsep=3pt]
  \item Many hyperparameters to tune (learning rate, depth, number of trees)
  \item Can overfit noisy data; sequential training is slower
\end{itemize}

\tpart{Applications}
\begin{itemize}[itemsep=2pt,topsep=3pt]
  \item Web search ranking, click-through-rate prediction
  \item Credit scoring and sales forecasting
\end{itemize}
\end{theory}

\begin{center}
\small
\begin{tabular}{@{}p{3.1cm}p{5.5cm}p{5.5cm}@{}}
\toprule
 & \textbf{AdaBoost} & \textbf{Gradient boosting}\\
\midrule
How errors are handled & Re-weights misclassified instances & Fits new learner to the residuals\\
Loss function & Exponential loss & Any differentiable loss\\
Combination & Weighted vote $\sum\alpha_th_t$ & Sum $F_0+\eta\sum h_m$\\
Weak learner & Usually a stump & Small tree (depth 3--8)\\
\bottomrule
\end{tabular}
\end{center}

%====================================================================
\section{Stacking}

Voting uses a fixed rule (majority, average). \textbf{Stacking} (stacked generalisation) asks: \emph{why not train a model to learn the best way to combine the predictions?}

\begin{definition}[: Stacking]
\textbf{Stacking} trains several different \textbf{base learners} (level-0) and then trains a \textbf{meta-learner} (level-1, also called the \emph{blender}) whose \emph{inputs are the predictions of the base learners} and whose output is the final prediction.
\end{definition}

\begin{figure}[!htbp]
\centering
\resizebox{\linewidth}{!}{%
\begin{tikzpicture}[
  b/.style={draw,rounded corners,thick,align=center,font=\small,minimum height=0.9cm},
  arr/.style={-{Stealth},thick}]
\node[b,fill=mlblue!12,draw=mlblue,minimum width=2.4cm,minimum height=1.4cm] (x) at (0,0) {Input\\$\vect x$};
\node[b,fill=gray!12,minimum width=2.8cm] (b1) at (4,1.7) {Base learner 1\\(k-NN)};
\node[b,fill=gray!12,minimum width=2.8cm] (b2) at (4,0) {Base learner 2\\(SVM)};
\node[b,fill=gray!12,minimum width=2.8cm] (b3) at (4,-1.7) {Base learner 3\\(random forest)};
\node[b,fill=mlorange!14,draw=mlorange,minimum width=2.2cm] (f) at (8.3,0) {Predictions\\$(\hat p_1,\hat p_2,\hat p_3)$};
\node[b,fill=mlpurple!14,draw=mlpurple,minimum width=2.8cm,minimum height=1.4cm] (m) at (12.1,0) {Meta-learner\\(logistic regression)};
\node[b,fill=mlgreen!18,draw=mlgreen] (y) at (15.6,0) {Final\\prediction};
\foreach \b in {b1,b2,b3} {\draw[arr] (x)--(\b); \draw[arr] (\b)--(f);}
\draw[arr] (f)--(m); \draw[arr] (m)--(y);
\node[font=\footnotesize,mlblue!70!black] at (4,2.6) {Level 0};
\node[font=\footnotesize,mlpurple] at (12.1,1.2) {Level 1};
\end{tikzpicture}}
\caption{Architecture of a stacking ensemble.}\label{fig:stack}
\end{figure}

\begin{explain}
The input goes to three \textbf{level-0} base learners (grey) that are different algorithms. Their three outputs are collected into a small table of predictions (orange), which becomes the \emph{input features} of the \textbf{level-1 meta-learner} (purple). The meta-learner has \emph{learned} how much to trust each base learner: for example it may learn that the SVM is most reliable and give it the largest weight. Its output is the final prediction (green).
\end{explain}

\paragraph{Avoiding leakage.} If the base learners predict on the \emph{same data they were trained on}, their predictions are too good and the meta-learner learns to over-trust them. The standard fix is to use \textbf{out-of-fold predictions}: for each training instance use the prediction of a base model that was trained \emph{without} that instance (cross-validation). A simpler variant, \textbf{blending}, holds out a separate validation set to train the meta-learner.

\begin{example}[: A meta-learner in action]
Three trained base models predict the probability that a customer will leave (churn) for one customer: $(0.80,\ 0.30,\ 0.60)$. A linear meta-learner has learned the weights $(0.5,\ 0.2,\ 0.3)$.
\[
  \hat p=0.5(0.80)+0.2(0.30)+0.3(0.60)=0.40+0.06+0.18=\mathbf{0.64}.
\]
Since $0.64\ge0.5$ the customer is predicted to \textbf{churn}. A simple average of the three would give $0.567$; stacking gave more influence to the first model because the meta-learner found it most reliable.
\end{example}

\begin{theory}{Stacking}
\tpart{Definition}
\textbf{Stacking} (stacked generalisation) is an ensemble technique in which the predictions of several diverse \emph{base learners} are used as inputs to a \emph{meta-learner} that is trained to produce the final prediction.

\tpart{Working}
\begin{enumerate}[itemsep=2pt,topsep=3pt]
  \item Split the training data into folds.
  \item Train each base learner and obtain \emph{out-of-fold predictions} for every training instance.
  \item Use these predictions as features and the true labels as targets to train the meta-learner.
  \item Retrain the base learners on all training data.
  \item For a new instance: get the base predictions, then pass them to the meta-learner.
\end{enumerate}

\tpart{Assumptions}
\begin{itemize}[itemsep=2pt,topsep=3pt]
  \item Base learners are diverse and make different errors
  \item Meta-learner is simple enough not to overfit
\end{itemize}

\tpart{Advantages}
\begin{itemize}[itemsep=2pt,topsep=3pt]
  \item Learns the best combination automatically
  \item Often more accurate than any single model or simple voting
\end{itemize}

\tpart{Disadvantages}
\begin{itemize}[itemsep=2pt,topsep=3pt]
  \item Complex, slow to train and tune
  \item Risk of information leakage and overfitting
  \item Hard to interpret
\end{itemize}

\tpart{Applications}
\begin{itemize}[itemsep=2pt,topsep=3pt]
  \item Machine-learning competitions (e.g.\ the Netflix Prize)
  \item High-stakes prediction where small accuracy gains matter
\end{itemize}
\end{theory}

%====================================================================
\section{Comparison and Practical Advice}

\begin{center}
\small
\begin{tabular}{@{}p{2.6cm}p{3.5cm}p{3.5cm}p{3.5cm}@{}}
\toprule
 & \textbf{Bagging / Random Forest} & \textbf{Boosting} & \textbf{Stacking}\\
\midrule
Training & Parallel & Sequential & Two levels\\
Data used & Bootstrap samples & Re-weighted data / residuals & Out-of-fold predictions\\
Main effect & Reduces \textbf{variance} & Reduces \textbf{bias} (and some variance) & Learns best combination\\
Base learner & Strong, high variance (deep trees) & Weak, high bias (stumps, shallow trees) & Diverse algorithms\\
Combination & Vote / average & Weighted sum & Meta-model\\
Overfitting risk & Low & Higher (noisy data) & Moderate\\
Parallelisable? & Yes & No & Partly\\
Examples & Random Forest, Extra-Trees & AdaBoost, XGBoost, LightGBM & StackingClassifier\\
\bottomrule
\end{tabular}
\end{center}

\begin{tryit}
Compare single models and ensembles on the \texttt{moons} dataset (two interleaving half circles, with noise).

\begin{lstlisting}
from sklearn.datasets import make_moons
from sklearn.model_selection import train_test_split
from sklearn.ensemble import (VotingClassifier, BaggingClassifier,
        RandomForestClassifier, AdaBoostClassifier,
        GradientBoostingClassifier, StackingClassifier)
from sklearn.linear_model import LogisticRegression
from sklearn.tree import DecisionTreeClassifier
from sklearn.svm import SVC

X, y = make_moons(n_samples=500, noise=0.30, random_state=42)
Xtr, Xte, ytr, yte = train_test_split(X, y, test_size=0.3,
                                      random_state=42)

lr = LogisticRegression()
rf = RandomForestClassifier(n_estimators=100, random_state=42)
sv = SVC(probability=True, random_state=42)

# soft voting
vote = VotingClassifier([("lr", lr), ("rf", rf), ("svc", sv)],
                        voting="soft").fit(Xtr, ytr)
print("soft voting:", vote.score(Xte, yte))

# bagging of 200 trees with out-of-bag score
bag = BaggingClassifier(DecisionTreeClassifier(), n_estimators=200,
        bootstrap=True, oob_score=True, random_state=42).fit(Xtr, ytr)
print("bagging:", bag.score(Xte, yte), "OOB:", bag.oob_score_)
\end{lstlisting}

\textbf{Results} on the 150 test points:

\medskip
\begin{center}
\begin{tabular}{@{}lc@{\qquad}lc@{}}
\toprule
\textbf{Model} & \textbf{Test accuracy} & \textbf{Model} & \textbf{Test accuracy}\\
\midrule
Single decision tree & 0.853 (train: 1.000) & AdaBoost (200 stumps) & 0.880\\
Logistic regression & 0.860 & Gradient boosting & 0.907\\
SVM & 0.907 & Bagging (200 trees) & 0.900 (OOB 0.891)\\
Hard voting & 0.913 & Pasting (200 trees) & 0.900\\
\textbf{Soft voting} & \textbf{0.933} & Random Forest & 0.900 (OOB 0.897)\\
Stacking & 0.900 & & \\
\bottomrule
\end{tabular}
\end{center}

\medskip
\textbf{What to notice:} the single tree \emph{memorises} the training set (train accuracy $1.0$) but scores only $0.853$ on test data; \textbf{bagging the same trees raises test accuracy to $0.90$}; and the OOB score ($0.891$, $0.897$) is close to the test accuracy without using the test set. \emph{Honest note:} with only $150$ test points one point is $0.67\%$, so differences of $1$--$2\%$ between ensembles are within noise; we should not claim that one ensemble is always the best.
\end{tryit}

\begin{warning}
\textbf{Common mistakes with ensembles.}
\begin{itemize}
  \item Combining models that are all \emph{the same} (no diversity) gives no gain.
  \item Evaluating a stacking model with base predictions made on the training data causes leakage.
  \item Using boosting on very noisy labels: the algorithm keeps chasing the noise.
  \item Assuming more trees always overfit: for Random Forests, more trees never hurt much, but for boosting too many stages can.
\end{itemize}
\end{warning}

%====================================================================
\begin{summary}
\begin{itemize}
  \item \textbf{Ensemble learning} combines several diverse models; it works because errors partly cancel (three $70\%$ classifiers give $78.4\%$ by majority vote, five give $83.7\%$).
  \item \textbf{Voting:} hard (majority of labels) or soft (average of probabilities, usually better).
  \item \textbf{Bagging} (with replacement) and \textbf{pasting} (without) train the same learner on random subsets; $\Var=\rho\sigma^2+\frac{1-\rho}{B}\sigma^2$; OOB $\approx36.8\%$ gives free validation.
  \item \textbf{Random Forest} = bagged trees + random feature subset at each split; gives feature importances.
  \item \textbf{AdaBoost:} re-weight misclassified instances; $\alpha_t=\frac12\ln\frac{1-\varepsilon_t}{\varepsilon_t}$; final $H=\operatorname{sign}\sum\alpha_th_t$.
  \item \textbf{Gradient boosting:} fit each tree to the current residuals; $F_m=F_{m-1}+\eta h_m$.
  \item \textbf{Stacking:} a meta-learner learns to combine base predictions; use out-of-fold predictions.
  \item Bagging reduces \textbf{variance}; boosting reduces \textbf{bias}; stacking learns the combination.
\end{itemize}
\end{summary}

%====================================================================
\section*{Exam Corner}
\addcontentsline{toc}{section}{Exam Corner}

\subsection*{A.\ Very Short Answer Questions (2 marks each)}
\begin{enumerate}
  \item Define ensemble learning. \qmarks{2}
  \item What is a weak learner? \qmarks{2}
  \item Differentiate between hard voting and soft voting. \qmarks{2}
  \item Differentiate between bagging and pasting. \qmarks{2}
  \item What is a bootstrap sample? \qmarks{2}
  \item What is an out-of-bag instance? What fraction of the data is OOB? \qmarks{2}
  \item How does a Random Forest differ from bagging of decision trees? \qmarks{2}
  \item What is feature importance in a Random Forest? \qmarks{2}
  \item State the formula for the weight $\alpha_t$ of a learner in AdaBoost. \qmarks{2}
  \item What does gradient boosting fit at each stage? \qmarks{2}
  \item What is a meta-learner? \qmarks{2}
  \item Which ensemble technique reduces variance and which reduces bias? \qmarks{2}
\end{enumerate}

\begin{modelans}[: Q6, Q9 and Q12]
\textbf{Q6.} When a bootstrap sample of size $n$ is drawn with replacement, some instances are never selected; they are the \emph{out-of-bag} instances of that sample. The chance that an instance is left out is $(1-1/n)^n\to e^{-1}\approx0.368$, so about \textbf{36.8\%} of the data is OOB for each model.\par
\textbf{Q9.} $\alpha_t=\dfrac12\ln\dfrac{1-\varepsilon_t}{\varepsilon_t}$, where $\varepsilon_t$ is the weighted error of the $t$-th weak learner.\par
\textbf{Q12.} Bagging (and Random Forests) mainly reduce \emph{variance}; boosting mainly reduces \emph{bias}.
\end{modelans}

\subsection*{B.\ Short Answer Questions (5 marks each)}
\begin{enumerate}
  \item Explain why ensemble methods improve accuracy. What is the role of diversity? \qmarks{5}
  \item Explain hard and soft voting classifiers with an example. \qmarks{5}
  \item Explain bagging with a diagram. How does out-of-bag evaluation work? \qmarks{5}
  \item Differentiate between bagging and pasting. \qmarks{5}
  \item Explain the Random Forest algorithm. Why do we choose a random subset of features at each split? \qmarks{5}
  \item Explain the AdaBoost algorithm. \qmarks{5}
  \item Explain gradient boosting with a suitable example. \qmarks{5}
  \item Explain stacking with a neat diagram. \qmarks{5}
  \item Differentiate between bagging and boosting. \qmarks{5}
\end{enumerate}

\begin{modelans}[: B9 (comparison for 5 marks)]
\begin{center}\small
\begin{tabular}{@{}p{2.6cm}p{5.2cm}p{5.2cm}@{}}
\toprule
\textbf{Basis} & \textbf{Bagging} & \textbf{Boosting}\\
\midrule
Training & Parallel; models independent & Sequential; each model depends on the previous\\
Data & Bootstrap samples & Re-weighted instances / residuals\\
Goal & Reduce variance & Reduce bias\\
Base learner & Strong, unstable (deep trees) & Weak (stumps)\\
Combination & Equal-weight vote/average & Weighted vote/sum\\
Noise & Robust & Sensitive\\
Example & Random Forest & AdaBoost, XGBoost\\
\bottomrule
\end{tabular}
\end{center}
\end{modelans}

\subsection*{C.\ Long Answer Questions (10 marks each)}
\begin{enumerate}
  \item What is ensemble learning? Explain the different types of ensemble methods with diagrams. \qmarks{10}
  \item Explain bagging, pasting and Random Forests. Discuss why they reduce variance. \qmarks{10}
  \item Explain the AdaBoost algorithm in detail and work through a numerical example. \qmarks{10}
  \item Explain gradient boosting and show with a numerical example how the error decreases over stages. \qmarks{10}
  \item Explain voting classifiers (hard, soft, weighted) and stacking. Give examples. \qmarks{10}
  \item Compare bagging, boosting and stacking in terms of training, bias--variance effect, advantages and disadvantages. \qmarks{10}
\end{enumerate}

\begin{modelans}[: C3 (key steps of the numerical part)]
Take $5$ instances with equal weights $0.2$. A weak learner misclassifies one instance: $\varepsilon_1=0.2$.\par
$\alpha_1=\frac12\ln\frac{0.8}{0.2}=0.693$. Correct instances: $0.2\times e^{-0.693}=0.1$; misclassified: $0.2\times e^{0.693}=0.4$. Sum $=4(0.1)+0.4=0.8$.\par
Normalised weights: correct $=0.1/0.8=0.125$, misclassified $=0.4/0.8=0.5$.\par
Next learner concentrates on the misclassified instance. Final model $H(\vect x)=\operatorname{sign}\big(\sum\alpha_th_t(\vect x)\big)$.
\end{modelans}

\subsection*{D.\ Numerical Problems}
\begin{enumerate}
  \item Three independent classifiers each have accuracy $0.8$. Find the accuracy of the majority vote. \textit{(Answer: $3(0.8)^2(0.2)+0.8^3=0.384+0.512=0.896$.)}
  \item Three classifiers give $P(\text{class }1)=(0.9,\,0.45,\,0.40)$. Find the hard- and soft-voting predictions. \textit{(Answer: hard votes $1,0,0\Rightarrow$ class $0$; soft: $\bar P=0.583\Rightarrow$ class $1$.)}
  \item Using weights $(2,1,1)$ for the classifiers of problem~2, find the weighted soft-voting probability. \textit{(Answer: $\frac{2(0.9)+0.45+0.40}{4}=\frac{2.65}{4}=0.6625\Rightarrow$ class $1$.)}
  \item For $n=100$ compute the probability that a given instance is \emph{not} in a bootstrap sample. \textit{(Answer: $(0.99)^{100}=0.366$; so $63.4\%$ are in the sample.)}
  \item $B=20$ models each have variance $9$ and pairwise correlation $0.3$. Find the variance of their average and the limit as $B\to\infty$. \textit{(Answer: $0.3(9)+\frac{0.7(9)}{20}=2.7+0.315=3.015$; limit $=2.7$.)}
  \item Find $\alpha_t$ for $\varepsilon_t=0.1$, $0.3$ and $0.5$ in AdaBoost. \textit{(Answer: $\frac12\ln9=1.099$; $\frac12\ln\frac73=0.424$; $0$.)}
  \item In AdaBoost with $10$ instances (weights $0.1$) a learner misclassifies $2$. Find $\varepsilon$, $\alpha$ and the new normalised weights. \textit{(Answer: $\varepsilon=0.2$, $\alpha=0.693$; correct $0.05\to0.0625$ each, wrong $0.2\to0.25$ each.)}
  \item Gradient boosting: $y=(2,4,9,11)$ at $x=1,2,3,4$. $F_0=\bar y$; a stump splits at $x\le2.5$; $\eta=0.1$. Find $F_1$ and the MSE before and after. \textit{(Answer: $F_0=6.5$; residuals $(-4.5,-2.5,2.5,4.5)$; stump $=(-3.5,3.5)$; $F_1=(6.15,6.15,6.85,6.85)$; MSE $13.25\to10.92$.)}
  \item A stacking meta-learner has weights $(0.5,0.2,0.3)$ and base predictions $(0.8,0.3,0.6)$. Find the final probability. \textit{(Answer: $0.40+0.06+0.18=0.64$.)}
\end{enumerate}

\subsection*{E.\ Multiple-Choice Questions}
\begin{enumerate}
  \item Combining the predictions of several models is called: \\
  (a) regularisation \quad (b) ensemble learning \quad (c) scaling \quad (d) pruning
  \item In soft voting the final class is chosen using: \\
  (a) majority of labels \quad (b) average of predicted probabilities \quad (c) random choice \quad (d) the first classifier
  \item Bagging samples the training data: \\
  (a) without replacement \quad (b) with replacement \quad (c) only once \quad (d) never
  \item The fraction of instances that are OOB for large $n$ is about: \\
  (a) 10\% \quad (b) 25\% \quad (c) 36.8\% \quad (d) 63.2\%
  \item Random Forests choose, at each split: \\
  (a) all features \quad (b) one fixed feature \quad (c) a random subset of features \quad (d) no feature
  \item Boosting mainly reduces: \\
  (a) variance \quad (b) bias \quad (c) data size \quad (d) noise
  \item Bagging mainly reduces: \\
  (a) bias \quad (b) variance \quad (c) irreducible error \quad (d) training time
  \item In AdaBoost, a learner with error $0.5$ gets weight $\alpha=$ \\
  (a) 1 \quad (b) $\ln2$ \quad (c) 0 \quad (d) $\infty$
  \item Gradient boosting fits each new tree to the: \\
  (a) original labels \quad (b) residuals \quad (c) random noise \quad (d) features
  \item The final model in stacking is called the: \\
  (a) base learner \quad (b) meta-learner \quad (c) weak learner \quad (d) stump
  \item Which method can be trained fully in parallel? \\
  (a) AdaBoost \quad (b) Gradient boosting \quad (c) Random Forest \quad (d) all of these
  \item Ensembles help most when members are: \\
  (a) identical \quad (b) accurate and diverse \quad (c) inaccurate \quad (d) correlated
\end{enumerate}
\textbf{Answers:} 1-b,\ 2-b,\ 3-b,\ 4-c,\ 5-c,\ 6-b,\ 7-b,\ 8-c,\ 9-b,\ 10-b,\ 11-c,\ 12-b.

\subsection*{F.\ Practice Problems}
\begin{enumerate}
  \item Compute by hand the accuracy of a majority vote of $7$ independent classifiers each $60\%$ accurate. \textit{(Answer: $0.710$, see Figure~\ref{fig:majority}.)}
  \item Continue the AdaBoost example: suppose a fourth learner has $\varepsilon_4=0.1$. Compute $\alpha_4$ and explain how it changes the combined vote.
  \item \textbf{Coding:} train a \texttt{RandomForestClassifier} on the \emph{wine} dataset (\texttt{load\_wine}), print the feature importances sorted, and compare accuracy with a single \texttt{DecisionTreeClassifier}.
  \item \textbf{Coding:} plot test accuracy of \texttt{GradientBoostingClassifier} against \texttt{n\_estimators} ($1$ to $300$) for learning rates $0.5$ and $0.05$; find where it starts to overfit.
  \item Draw a diagram of a stacking ensemble and label base learners, predictions and meta-learner.
\end{enumerate}

\begin{revision}
\begin{center}\small
\begin{tabular}{@{}p{3.6cm}p{9.6cm}@{}}
\toprule
\textbf{Topic} & \textbf{One-line recall}\\
\midrule
Ensemble & Many diverse models combined; better and more stable than one.\\
Majority vote & $P=\sum_{k>n/2}\binom nk p^k(1-p)^{n-k}$; 3 models of 70\% $\to$ 78.4\%.\\
Hard / soft voting & Majority of labels / average of probabilities (soft often better).\\
Bagging / pasting & Bootstrap (with replacement) / without replacement; vote or average.\\
OOB & $\approx36.8\%$ left out of each bootstrap sample; free validation.\\
Variance of average & $\rho\sigma^2+\frac{1-\rho}{B}\sigma^2$.\\
Random Forest & Bagged trees $+$ random features per split ($\sqrt p$); importances.\\
AdaBoost & $\alpha=\frac12\ln\frac{1-\varepsilon}{\varepsilon}$; $w\leftarrow w e^{\mp\alpha}$; $H=\operatorname{sign}\sum\alpha h$.\\
Gradient boosting & Fit residuals; $F_m=F_{m-1}+\eta h_m$.\\
Stacking & Meta-learner on out-of-fold base predictions.\\
Bias--variance & Bagging: variance $\downarrow$; boosting: bias $\downarrow$.\\
\bottomrule
\end{tabular}
\end{center}
\end{revision}
